% ===== §2 The Conceptual Archaeology of Luxury =====
\section{The Conceptual Archaeology of Luxury}
\label{sec:archaeology}

\noindent
Before the three doctrines can be developed, the concept of luxury itself must be excavated---its historical sedimentation exposed, its internal tensions identified, its presuppositions made visible. The concept of luxury is not a neutral descriptive category; it is a historically produced and politically contested one, whose apparent self-evidence conceals a series of structural assumptions that the relational account will need to challenge.

\subsection{Historical Forms: From Royal Prerogative to Democratic Consumption}
\label{subsec:history}

\noindent
The history of luxury is, in large part, a history of exclusion and its gradual relaxation. In pre-modern societies, luxury was constitutively tied to social position: the right to consume certain materials, wear certain garments, or display certain objects was legally reserved for specific social ranks, enforced by sumptuary laws that prescribed what each rank could and could not wear, eat, or own. The luxury object was, in this period, primarily a \emb{marker of social position}---not merely an indicator of wealth but a legally constituted symbol of rank. The diamond was not merely expensive; it was, in many jurisdictions, legally restricted to those of sufficiently elevated social station.

The gradual dismantling of sumptuary law across the early modern period transformed the social function of luxury without eliminating its positional character. As legal restriction gave way to market restriction---luxury becoming available in principle to anyone who could afford it, rather than only to those of appropriate rank---the positional character of luxury was reproduced through a different mechanism: the price mechanism. The diamond ring remained a marker of social distinction, but the distinction it marked shifted from legal rank to economic class. Thorstein Veblen's analysis of \emb{conspicuous consumption} \citep{veblen1899} captures the structure of this transition: luxury consumption in the market era functions as a display of surplus---the ability to spend beyond necessity---that signals to others the consumer's position in the economic hierarchy.

The twentieth century brought a further transformation: the \emb{democratisation of luxury}, in which the luxury industry developed strategies for extending its market to the aspirational middle class without losing the exclusivity that constituted its appeal. The development of accessible luxury lines, licensed products, and the expansion of the luxury market into new territories (cosmetics, perfume, fashion accessories) allowed the luxury industry to maintain the symbolic apparatus of exclusivity---the brand heritage, the artisanal craftsmanship, the association with high social status---while dramatically expanding the consumer base. The result is the contemporary luxury landscape: a market in which luxury objects are available to a far wider range of consumers than at any previous point in history, while remaining organised by the logic of social distinction that has structured luxury since the era of sumptuary law.

\subsection{Scarcity, Craft, and Beauty: The Three Traditional Dimensions}
\label{subsec:traditional}

\noindent
The traditional definition of luxury organises itself around three dimensions that correspond to three different sources of value.

\emb{Scarcity} is the most fundamental: the luxury object is, by definition, not available to everyone. This scarcity may be natural (the rarity of the gemstone, the difficulty of sourcing the material), artificial (the deliberate restriction of production to maintain exclusivity), or temporal (the limited edition that will not be reproduced). In all cases, the scarcity of the luxury object is constitutive of its value: what everyone can have is, by that very fact, not luxurious. The luxury object's value is relational---it is defined by its position in a system of differential access, not by any intrinsic property of the object itself.

\emb{Craft} is the second dimension: the luxury object is distinguished from the merely expensive object by the quality of its making. The handmade watch, the hand-stitched garment, the individually cut gemstone represent forms of human labour that are increasingly rare in an age of industrial production, and their rarity gives them a value that goes beyond the material cost of their production. Craft in luxury objects is not merely a productive process but a \emb{form of attention}: the craftsperson who spends weeks on a single watch movement, or months on a single garment, is performing a kind of attentive care for the object that is qualitatively different from the standardised repetition of industrial production. This attentive care is what the luxury object carries within itself: the trace of a human being's sustained engagement with the material.

\emb{Beauty} is the third dimension: the luxury object is aesthetically distinguished from ordinary objects. This aesthetic distinction is not merely a matter of decoration---it reflects a philosophy of design in which the object's form is the expression of a vision of the beautiful, developed through the traditions of a particular craft culture and the aesthetic judgements of particular designers. The beauty of a Bulgari jewel is not the accidental result of expensive materials; it is the product of a long tradition of aesthetic thought about how natural materials can be given form that honours both their natural properties and the human capacity for aesthetic judgement.

These three dimensions interact: scarcity without craft or beauty produces merely the expensive; craft without scarcity or beauty produces the skilled but ordinary; beauty without scarcity or craft produces the aesthetically pleasing but not luxurious. The canonical luxury object---the hand-cut diamond in a hand-crafted setting, designed by a master of the tradition---combines all three.

\subsection{The Instability of the Luxury/Necessity Distinction}
\label{subsec:instability}

\noindent
The distinction between luxury and necessity that seems to anchor the concept of luxury is, on examination, unstable in ways that are philosophically significant. What counts as a necessity and what counts as a luxury is not a natural fact but a historically and culturally variable social construction---and the variation reveals the political character of the distinction.

At one level, the instability is obvious: what was a luxury in one era is a necessity in another (indoor plumbing, the telephone, electric lighting), and what is a luxury in one society is a necessity in another (the car in the suburban United States, the mobile phone in contemporary rural Africa). The relativity of the luxury/necessity distinction to historical and cultural context means that the concept of luxury is not defined by any intrinsic property of the objects classified as luxuries, but by their position in a particular social and economic context.

At a deeper level, the instability is philosophically revealing. The luxury/necessity distinction presupposes a conception of the human being and human needs that is itself contested: a conception of the \emb{bare minimum}---the floor below which life is not sustainable---against which surplus is defined. But this floor is not naturally given; it is socially constructed, and the construction reflects relations of power. What is counted as a necessity is what is deemed necessary for the reproduction of labour power; what is counted as a luxury is what exceeds this minimum. The political economy of the luxury/necessity distinction is therefore not neutral: it reflects and reproduces the interests of those who benefit from keeping the definition of necessity narrow.

This instability has a direct implication for the relational account of luxury developed in this paper: if the luxury/necessity distinction is already unstable within the dominant framework, the relational framework's proposal---that what constitutes true luxury is the degree of coupling between an object and a relational subject's spatiotemporal structure---represents not an arbitrary redefinition but a principled attempt to locate the value of luxury in a dimension that the dominant framework has systematically missed.

\subsection{The Contemporary Luxury Industry: Political Economy}
\label{subsec:industry}

\noindent
The contemporary luxury industry is one of the most economically significant sectors of the global economy, and its political economy reveals the tensions within the concept of luxury that the three doctrines of this paper will develop. Several features of the contemporary luxury industry are particularly relevant.

\emb{The brand as the primary unit of value}: In the contemporary luxury market, the primary unit of value is not the individual object but the brand---the accumulated cultural capital of a name, a heritage, and a set of aesthetic associations that gives any object bearing the brand's mark a premium over its unbranded equivalent. The brand is a symbolic entity: it exists in the collective imagination of consumers and non-consumers alike, and its value depends on the maintenance of its symbolic associations. Bulgari, Cartier, Tiffany, Van Cleef \& Arpels: these names carry a symbolic weight that is the product of decades of marketing, cultural positioning, and the careful management of the brand's associations with beauty, excellence, and social distinction.

The primacy of the brand in contemporary luxury has a direct implication for the Baudrillard analysis developed in \xref{sec:baudrillard}: the luxury object's value is organised primarily by its symbolic position rather than by its material properties or its craft quality. The diamond in a Tiffany setting is valued not primarily for the diamond's natural properties or the setting's craft quality, but for the Tiffany name---for what the Tiffany name means within the symbolic order of consumer society.

\emb{The globalisation of luxury and the management of aspiration}: The contemporary luxury industry has globalised its market while managing the aspiration that drives consumption by maintaining the appearance of exclusivity even as the consumer base expands. The luxury industry's characteristic strategy is the creation of \emb{accessible luxury}---products priced at the lower end of the luxury range, which allow aspirational consumers to participate in the symbolic system of luxury without access to its highest tiers. The handbag that costs a month's salary rather than a year's performs the same symbolic function as the handbag that costs a year's salary, for consumers at different points in the economic hierarchy.

\emb{The tension between craft and scale}: The contemporary luxury industry faces a structural tension between its dependence on craft as a source of value and its need for scale as a commercial enterprise. The handmade watch that took a craftsperson weeks to produce cannot be manufactured at the scale required by a publicly traded corporation; the solution, characteristic of the industry, is the combination of artisanal quality at the level of finishing and detail with industrial production at the level of components and structure. This tension reveals the gap between the craft mythology of luxury and its industrial reality---a gap that is managed through the careful control of the narrative of luxury, rather than through the actual maintenance of artisanal production at scale.

\emb{The digital transformation of luxury}: The emergence of digital luxury---NFTs, virtual luxury goods, the use of luxury brands in digital environments---represents a further development of the sign exchange value logic identified by Baudrillard: the luxury object is increasingly abstracted from any material substrate, becoming a pure sign without a referent. The digital luxury object is, in Baudrillard's terms, a \emb{simulacrum}---a sign that refers only to other signs, with no connection to any material reality \citep{baudrillard1981}. The implications of this development for the relational account of luxury are direct: a digital luxury object, which has no material substrate in which a subject's spatiotemporal structure can be embedded, cannot in principle be a relational luxury object. The digital transformation of luxury is therefore the fullest expression of the tension between the sign exchange value logic of contemporary luxury and the relational reality of intimate life.

\begin{claim}[The conceptual archaeology's conclusion]
The concept of luxury is not a neutral descriptive category but a historically produced, politically contested, and structurally unstable one, organised in its dominant forms around the dimensions of scarcity, craft, and beauty, and in its contemporary market form around the primacy of brand as a symbolic entity. The relational account's proposal---that the true measure of luxury in intimate relations is the degree of coupling between an object and a relational subject's spatiotemporal structure---is not a departure from the concept of luxury but a recovery of its most fundamental dimension: the dimension of attentive care, of the particular, of the irreplaceable.
\end{claim}
